% !TEX root = ../main.tex
%% Corrigés — Chapitre 18 — Sommes de variables aléatoires
%% Tous les résultats ont été vérifiés par programme (calcul exact et énumération des issues).

\corrigesChapitre{18}{Sommes de variables aléatoires}

\begin{deuxcolonnes}

\begin{corrige}{1}{Espérance et variance : rappels}
\begin{questions}
  \item $E(X)=-0{,}2+0+1{,}2+1=\res{2}$ ;\\ $E\left(X^{2}\right)=0{,}4+0+3{,}6+5=9$ ;\\ $V(X)=9-2^{2}=\res{5}$.
  \item $E(Y)=3\times2-1=\res{5}$ et $V(Y)=9\times5=\res{45}$ ;\\ $E(Z)=-4+4=\res{0}$ et $V(Z)=4\times5=\res{20}$.
\end{questions}
\end{corrige}

\begin{corrige}{2}{Linéarité de l'espérance}
\methode{$E(X+Y)=E(X)+E(Y)$}
\begin{questions}
  \item $3{,}5+3{,}5=\res{7}$
  \item $\frac12+\frac12=\res{1}$
  \item $S=X_1+X_2$, où $X_i$ vaut $1$ si la $i$-ème boule est rouge : $E(S)=2\times\frac58=\res{\frac54}$.
  \item $3\times3{,}5+2=\res{12{,}5}$
  \item $10\times\frac16+2\times\frac56=\res{\frac{10}{3}}$ points
\end{questions}
\end{corrige}

\begin{corrige}{3}{Additivité de la variance}
\methode{variables indépendantes}
\begin{questions}
  \item $\frac{35}{12}+\frac{35}{12}=\res{\frac{35}{6}}$
  \item $\frac14+\frac14=\res{\frac12}$
  \item $2\times0{,}4\times0{,}6=\res{\frac{12}{25}}$
  \item $3^{2}\times\frac{35}{12}=\res{\frac{105}{4}}$
  \item $E(X)=p$ et $E\left(X^{2}\right)=p$ (car $X^{2}=X$), donc $V(X)=p-p^{2}=\res{p(1-p)}$.
\end{questions}
\end{corrige}

\begin{corrige}{4}{La somme de deux dés}
\begin{questions}
  \item Les $36$ issues sont équiprobables :\\ $P(S=k)=\frac{6-|k-7|}{36}$ pour $2\leq k\leq12$, soit $\frac1{36}$, $\frac2{36}$, \dots,
    $\frac6{36}$ (pour $k=7$), \dots, $\frac1{36}$.
  \item $E(S)=\res{7}$ et $E\left(S^{2}\right)=\frac{1\,974}{36}$, d'où $V(S)=\frac{1\,974}{36}-49=\res{\frac{35}{6}}$.
  \item $E(S)=\frac72+\frac72=7$ et, les dés étant indépendants, $V(S)=2\times\frac{35}{12}=\frac{35}{6}$.
  \item $E(D)=\frac72-\frac72=\res{0}$ et $V(D)=V(R)+V(B)=\res{\frac{35}{6}}$ : les variances s'ajoutent.
\end{questions}
\end{corrige}

\end{deuxcolonnes}

%% ---------------------------------------------------------------------
\begin{corrige}{5}{Applications}
\begin{questions}
  \item $X\sim\mathcal{B}\left(10\pv\frac12\right)$ et $X=X_1+\dots+X_{10}$, où $X_i$ vaut $1$ si le $i$-ème lancer donne pile : $E(X)=\res{5}$
    et $V(X)=10\times\frac14=\res{\frac52}$.
  \item a) $X\sim\mathcal{B}\left(5\pv\frac38\right)$ et $Y\sim\mathcal{B}\left(5\pv\frac12\right)$ : $E(X)=\res{\frac{15}{8}}$, $V(X)=\res{\frac{75}{64}}$,
    $E(Y)=\res{\frac52}$, $V(Y)=\res{\frac54}$.\\
    b) $E(S)=\frac{15}{8}+\frac52=\res{\frac{35}{8}}$. $S$ compte les boules non vertes : $S\sim\mathcal{B}\left(5\pv\frac78\right)$, d'où
    $V(S)=5\times\frac78\times\frac18=\res{\frac{35}{64}}$.\\
    c) $V(X)+V(Y)=\frac{155}{64}\neq V(S)$ : $X$ et $Y$ ne sont pas indépendantes (plus il y a de rouges, moins il peut y avoir de bleues),
    donc les variances ne s'ajoutent pas.
  \item $E(S)=\res{10{,}5}$ et $V(S)=3\times\frac{35}{12}=\res{\frac{35}{4}}$ ; $E(M)=\res{3{,}5}$ et $V(M)=\frac19V(S)=\res{\frac{35}{36}}$.
  \item Le gain d'un lancer vaut $5$ ou $-2$ avec la probabilité $\frac12$ : espérance $1{,}5$, variance $\frac{25+4}{2}-1{,}5^{2}=12{,}25$.
    Pour $4$ lancers indépendants : $E(G)=\res{6}$ et $V(G)=\res{49}$.
  \item $E(X)=\res{15}$ et $V(X)=\res{6}$ ; $E(F)=\frac{15}{25}=\res{0{,}6}$ et $V(F)=\frac{6}{25^{2}}=\res{\frac{6}{625}}$.
\end{questions}
\end{corrige}

%% ---------------------------------------------------------------------
\begin{corrige}{6}{Moyenne de plusieurs mesures}
\begin{questions}
  \item $E(M_n)=\res{\mu}$ et $\sigma(M_n)=\res{\frac{0{,}5}{\sqrt n}}$ : la moyenne est centrée sur la vraie valeur et elle est
    d'autant moins dispersée que $n$ est grand.
  \item $\frac{0{,}5}{\sqrt n}\leq0{,}1\iff\sqrt n\geq5\iff\res{n\geq25}$.
  \item $\sigma(S_{25})=\sqrt{25}\times0{,}5=\res{2{,}5\text{ mm}}$, alors que $\sigma(M_{25})=\res{0{,}1\text{ mm}}$ : la somme se disperse,
    la moyenne se concentre.
\end{questions}
\end{corrige}

%% ---------------------------------------------------------------------
\begin{corrige}{7}{Situations complexes}
\begin{questions}
  \item $X\sim\mathcal{B}(100\pv0{,}05)$ : $E(X)=\res{5}$ et $V(X)=\res{4{,}75}$ ; $E(F)=\res{0{,}05}$ et $\sigma(F)=\frac{\sqrt{4{,}75}}{100}\approx\res{0{,}022}$.
    En moyenne, 5\,\% des composants d'un lot sont défectueux, avec des écarts typiques d'environ $2$ points.
  \item $E(G)=20\times\frac{7}{216}+10\times\frac{707}{2\,592}=\frac{1\,680+7\,070}{2\,592}\approx\res{3{,}38}$, puis $E(N)\approx3{,}38-8=\res{-4{,}62}$ :
    le jeu est défavorable au joueur.
  \item $X\sim\mathcal{B}(50\pv0{,}7)$ : $E(X)=35$, $V(X)=10{,}5$ ; $Y\sim\mathcal{B}(50\pv0{,}3)$ : $E(Y)=15$, $V(Y)=10{,}5$. Mais $X+Y=50$ est
    constante : $E(X+Y)=\res{50}$ et $V(X+Y)=\res{0}\neq V(X)+V(Y)$, car $X$ et $Y$ ne sont pas indépendantes ($Y=50-X$).
  \item $X\sim\mathcal{B}(200\pv0{,}6)$ : $E(X)=\res{120}$, $\sigma(X)=\sqrt{48}\approx\res{6{,}93}$ ; $E(F)=\res{0{,}6}$, $\sigma(F)\approx\res{0{,}034\,6}$ ;
    intervalle $\res{\intFF{0{,}531}{0{,}669}}$.
  \item $E(N)=3\times0{,}9=\res{2{,}7}$ et $V(N)=3\times0{,}09=\res{0{,}27}$ ; $N\sim\mathcal{B}(3\pv0{,}9)$, ce qui confirme ces valeurs ;
    $P(N\geq2)=3\times0{,}9^{2}\times0{,}1+0{,}9^{3}=\res{0{,}972}$.
\end{questions}
\end{corrige}

%% ---------------------------------------------------------------------
\begin{corrige}{8}{Devoir de synthèse}
\begin{questions}
  \item $3\times3{,}5=\res{10{,}5}$.
  \item $2\times4+3\times6-1=\res{25}$.
  \item Chaque lancer rapporte en moyenne $3\times\frac12-1\times\frac12=1$ point : $E(G)=\res{4}$.
  \item $V(2X-Y+1)=4V(X)+V(Y)=\res{23}$.
  \item $X\sim\mathcal{B}(20\pv0{,}4)$ : $E(X)=\res{8}$ et $V(X)=\res{4{,}8}$.
  \item $X\sim\mathcal{B}\left(15\pv\frac58\right)$ : $E(X)=\res{\frac{75}{8}}$ et $\sigma(X)=\sqrt{\frac{225}{64}}=\res{\frac{15}{8}}$.
  \item $E(S_{100})=\res{350}$ et $\sigma(S_{100})=10\sqrt{\frac{35}{12}}\approx\res{17{,}08}$ ; $E(M_{100})=\res{3{,}5}$ et $\sigma(M_{100})\approx\res{0{,}171}$.
  \item $E(F)=\res{0{,}45}$ et $V(F)=\frac{0{,}45\times0{,}55}{80}=\res{0{,}003\,093\,75}$ ; $\sigma(F)\approx0{,}055\,6$ et l'intervalle est $\res{\intFF{0{,}339}{0{,}561}}$.
  \item $N\sim\mathcal{B}(4\pv0{,}85)$ : $E(N)=\res{3{,}4}$, $V(N)=\res{0{,}51}$ ; $P(N\geq3)=4\times0{,}85^{3}\times0{,}15+0{,}85^{4}\approx\res{0{,}890}$.
\end{questions}
\end{corrige}

%% ---------------------------------------------------------------------
\begin{corrige}{9}{Un jeu répété}
\begin{questions}
  \item $P(R=10)=\frac16$, $P(R=2)=\frac13$, $P(R=0)=\frac12$. Donc $E(R)=\frac{10}{6}+\frac{2}{3}=\res{\frac73}$ ; $E\left(R^{2}\right)=\frac{100}{6}+\frac43=18$, et
    $V(R)=18-\frac{49}{9}=\res{\frac{113}{9}}\approx12{,}6$.
  \item $E(G)=\frac73-m=0$ pour $\res{m=\frac73\approx2{,}33\text{ €}}$.
  \item $E(G)=\frac73-3=\res{-\frac23}$ et $V(G)=V(R)=\res{\frac{113}{9}}$ (une constante ne change pas la variance).
  \item Par linéarité et indépendance : $E(T)=-\frac{200}{3}\approx\res{-66{,}7\text{ €}}$ et $\sigma(T)=\sqrt{100\times\frac{113}{9}}\approx\res{35{,}4\text{ €}}$.
    En moyenne, le joueur perd près de $67$\,€ en $100$ parties ; l'écart-type montre qu'un gain final reste possible, mais peu probable.
\end{questions}
\end{corrige}

%% ---------------------------------------------------------------------
\begin{corrige}{10}{Combien de faces différentes ?}
\begin{questions}
  \item Chaque face apparue compte pour $1$ : $\res{N=Y_1+Y_2+\dots+Y_6}$.
  \item $Y_j=0$ si la face $j$ n'apparaît à aucun des $6$ lancers indépendants : $P(Y_j=0)=\left(\frac56\right)^{6}\approx0{,}335$, et
    $E(Y_j)=P(Y_j=1)=\res{1-\left(\frac56\right)^{6}}$.
  \item Par linéarité : $E(N)=\res{6\left(1-\left(\frac56\right)^{6}\right)\approx3{,}99}$.
  \item $P(N=6)=\frac{6!}{6^{6}}\approx\res{0{,}015}$ (les six résultats forment une permutation). $P(Y_1=0\text{ et }Y_2=0)=\left(\frac46\right)^{6}\approx0{,}088$,
    alors que $P(Y_1=0)\times P(Y_2=0)=\left(\frac56\right)^{12}\approx0{,}112$ : les variables ne sont $\res{\text{pas indépendantes}}$. Le calcul reste
    valable, car la linéarité de l'espérance n'exige pas l'indépendance.
  \item Le même raisonnement donne $E(N)=6\left(1-\left(\frac56\right)^{n}\right)$. Puis $E(N)\geq5\iff\left(\frac56\right)^{n}\leq\frac16\iff
    n\geq\frac{\ln(1/6)}{\ln(5/6)}\approx9{,}8$ : à partir de $\res{10\text{ lancers}}$.
\end{questions}
\end{corrige}
